\section{A determinant obstruction from field norms}
\label{norm:section}

We construct random columns modulo a prime power, each carrying a label in
a finite field.  Three distinct labels will force the determinant residue
to have no small integer representative.  The construction first expresses
a field norm as a sum of determinants, then places these summands at one
specified coefficient in a base expansion.

\subsection{The norm polynomial}

Fix
\begin{equation}\label{norm:fixed-parameters}
 d=41,\qquad M=\binom{4d-1}{d},\qquad
 T=\binom{M}{3},\qquad k=T^2+1.
\end{equation}
Let $r$ be an odd prime.  We recall the elementary finite-field facts
needed here; see \cite[Chapter~2]{LidlNiederreiter1994} for background.
In a splitting field over $\mathbb F_r$, the roots of
$X^{r^d}-X$ form a field: the Frobenius identity shows closure under
addition, subtraction and multiplication, and nonzero roots are closed
under inversion.  The derivative is $-1$, so there are exactly $r^d$
distinct roots.  This field, denoted $K=\mathbb F_{r^d}$, consequently
has dimension $d$ over $\mathbb F_r$.  Its automorphism $t\mapsto t^r$
has $d$th power equal to the identity, and its fixed elements are exactly
the $r$ roots of $X^r-X$, namely $\mathbb F_r$.
For $t\in K$, the product $\prod_{j=0}^{d-1}t^{r^j}$ is therefore fixed
by Frobenius and belongs to $\mathbb F_r$; it is nonzero if $t\ne0$.
Choose an $\mathbb F_r$-basis $\beta_1,\ldots,\beta_d$ of $K$.
For a group $X=(X_{i\nu})_{1\le i\le3,\,1\le\nu\le d}$ of $3d$
variables, define
\[
 Z(X)=\left(\sum_{\nu=1}^d\beta_\nu X_{1\nu},
             \sum_{\nu=1}^d\beta_\nu X_{2\nu},
             \sum_{\nu=1}^d\beta_\nu X_{3\nu}\right)^{\mathsf T}.
\]
For three disjoint variable groups $X^{(1)},X^{(2)},X^{(3)}$, put
\[
 \mathcal D=\det\bigl(Z(X^{(1)}),Z(X^{(2)}),Z(X^{(3)})\bigr).
\]
Let $\sigma$ act on this polynomial ring by applying $c\mapsto c^r$
to coefficients in $K$ and fixing every variable.  Define
\begin{equation}\label{norm:polynomial}
 \mathcal N(X^{(1)},X^{(2)},X^{(3)})
   =\prod_{j=0}^{d-1}\sigma^j(\mathcal D).
\end{equation}
The action of $\sigma$ permutes the factors, so $\mathcal N$ has
coefficients in $\mathbb F_r$.  It is homogeneous of degree $d$ in each
group.  At arguments in $\mathbb F_r^{3d}$, it evaluates to the field norm
of the corresponding determinant:
\begin{align*}
 \mathcal N(U^{(1)},U^{(2)},U^{(3)})
  &=\operatorname{Nm}_{K/\mathbb F_r}
       \det\bigl(Z(U^{(1)}),Z(U^{(2)}),Z(U^{(3)})\bigr),\\
 \operatorname{Nm}_{K/\mathbb F_r}(t)&=\prod_{j=0}^{d-1}t^{r^j}.
\end{align*}
Indeed, coefficient conjugation commutes with evaluation at such arguments.
Transposing two variable groups changes the sign of each factor in
\eqref{norm:polynomial}.  Since $d$ is odd, $\mathcal N$ is alternating.

\begin{lemma}\label{norm:decomposition}
There exist homogeneous degree-$d$ polynomials $f_{ia}\in\mathbb F_r[X]$
(allowing the zero polynomial), for
$1\le i\le3$ and $0\le a<T$, such that
\begin{equation}\label{norm:determinant-sum}
 \mathcal N(X^{(1)},X^{(2)},X^{(3)})
   =\sum_{a=0}^{T-1}
      \det\bigl(f_{ia}(X^{(j)})\bigr)_{1\le i,j\le3}.
\end{equation}
For each $a$, the three polynomials are obtained from three distinct
degree-$d$ monomials by multiplying the first by a scalar that may be zero.
\end{lemma}

\begin{proof}
List the degree-$d$ monomials in $3d$ variables as $m_1,\ldots,m_M$.
Expand $\mathcal N$ in products
$m_p(X^{(1)})m_q(X^{(2)})m_s(X^{(3)})$.
If two of $p,q,s$ coincide, the coefficient equals its negative under the
corresponding transposition, and hence is zero because $r$ is odd.
For each $p<q<s$, let $c_{pqs}$ be the coefficient of
$m_p(X^{(1)})m_q(X^{(2)})m_s(X^{(3)})$.
Alternation identifies the other five coefficients with the signs in
\[
 c_{pqs}\det
 \begin{pmatrix}
  m_p(X^{(1)})&m_p(X^{(2)})&m_p(X^{(3)})\\
  m_q(X^{(1)})&m_q(X^{(2)})&m_q(X^{(3)})\\
  m_s(X^{(1)})&m_s(X^{(2)})&m_s(X^{(3)})
 \end{pmatrix}.
\]
Enumerate the $T$ triples $p<q<s$ by $a$, and take
\[
 (f_{1a},f_{2a},f_{3a})=(c_{pqs}m_p,m_q,m_s).
\]
This argument uses no division by $3!$, so it applies also when $r=3$.
\end{proof}

For a label $\xi\in K$, set $z(\xi)=(1,\xi,\xi^2)^{\mathsf T}$ and
let $X(\xi)\in\mathbb F_r^{3d}$ be its coordinates in the chosen basis,
so that $Z(X(\xi))=z(\xi)$.
The Vandermonde determinant gives
\begin{equation}\label{norm:distinct-labels}
 \mathcal N(X(\xi_1),X(\xi_2),X(\xi_3))
  =\operatorname{Nm}_{K/\mathbb F_r}
        \left(\prod_{1\le i<j\le3}(\xi_j-\xi_i)\right)\ne0
\end{equation}
whenever the three labels are distinct.
Although the basis and the coefficients $f_{ia}$ may depend on $r$, the
numbers $d,M,T,k$ in \eqref{norm:fixed-parameters} do not.

\subsection{Encoding the polynomial in a base expansion}

Two prime choices will keep the construction parameters polynomial in
$r$. We use the following sufficiently-large form of Bertrand's postulate;
Appendix~\ref{app:prime-interval} gives its elementary binomial proof,
following Erd\H{o}s~\cite{Erdos1932}.

\begin{lemma}\label{norm:prime-interval}
For every sufficiently large integer $n$, there is a prime $p$ with
$n<p\le2n$.
\end{lemma}

Set $L=r^{10}$.  For all sufficiently large $r$, apply
Lemma~\ref{norm:prime-interval} with $n=100k^2L^3$ to choose a prime
$B$ satisfying
\begin{equation}\label{norm:base}
 100k^2L^3<B\le200k^2L^3,
 \qquad h=B^k,\qquad
 \tau=\left\lfloor\frac{B^{k-1}}2\right\rfloor,
 \qquad R_h=\mathbb Z/h\mathbb Z.
\end{equation}
For $0\le a<T$, designate positions in rows $1,2,3$, respectively, by
\begin{equation}\label{norm:positions}
 i_a=a,\qquad j_a=Ta,\qquad \ell_a=T^2-(T+1)a.
\end{equation}
These positions lie in $\{0,\ldots,k-1\}$ and are distinct within each
row: they lie respectively in the intervals $[0,T-1]$, $[0,T^2-T]$,
and $[1,T^2]$.  The only designated positions whose sum is $k-1$ are
the matched triples:
\begin{equation}\label{norm:matching}
 i_a+j_{a'}+\ell_{a''}=k-1
 \quad\Longleftrightarrow\quad a=a'=a''.
\end{equation}
To prove the forward implication, rewrite the equation as
$a+Ta'=(T+1)a''$.  Reduction modulo $T$, together with
$0\le a,a''<T$, gives $a=a''$; substitution then gives $a'=a''$.
The reverse implication follows from \eqref{norm:positions}.

A random column $c=(c_1,c_2,c_3)^{\mathsf T}\in R_h^3$ is now defined
as follows.  First choose $\xi$ uniformly from $K$.  For each row
$1\le i\le3$ and each position $0\le v<k$, choose an integer digit
$c_{iv}\in\{1,\ldots,L\}$.  Its required residue modulo $r$ is
$f_{1a}(X(\xi))$ when $(i,v)=(1,i_a)$,
$f_{2a}(X(\xi))$ when $(i,v)=(2,j_a)$, and
$f_{3a}(X(\xi))$ when $(i,v)=(3,\ell_a)$.
At every other position its required residue is zero.
Conditional on $\xi$, choose all $3k$ digits independently and uniformly
subject to these requirements, and set
\begin{equation}\label{norm:column}
 c_i=\sum_{v=0}^{k-1}c_{iv}B^v\pmod h.
\end{equation}
Every required residue has exactly $L/r=r^9$ possible digits, including
the zero residue.  In particular a digit prescribed to be zero modulo
$r$ is still a positive integer.  Since $L<B$, the lowest digit is a
unit modulo $B$; thus every entry of $c$ is a unit in $R_h$.
Different columns use independent labels and fresh independent digits,
even when their labels coincide.

\begin{lemma}[Determinant obstruction]\label{norm:obstruction}
Let $C$ be a matrix of three columns constructed by
\eqref{norm:column}, with labels $\xi_1,\xi_2,\xi_3$.
If these labels are distinct, then the residue $\det C\in R_h$ has no
integer representative in $[-\tau,\tau]$.
\end{lemma}

\begin{proof}
Write $c_{iv}^{(j)}$ for the digit at row $i$, position $v$, in column
$j$.  Form the integer polynomial
\[
 F(Y)=\det\left(\sum_{v=0}^{k-1}c_{iv}^{(j)}Y^v\right)_{1\le i,j\le3}
      =\sum_{u=0}^{3(k-1)}\gamma_uY^u.
\]
For a fixed exponent $u$ and each of the six determinant permutations,
choosing the first two digit positions determines the third.  Therefore
\begin{equation}\label{norm:coefficient-bound}
 |\gamma_u|\le A_0:=6k^2L^3.
\end{equation}
Alternatively expanding by rows, $\gamma_{k-1}$ is the sum of
determinants of digit rows at positions whose sum is $k-1$.
Modulo $r$, all terms with an undesignated position vanish.
By \eqref{norm:matching}, the surviving terms are exactly those in
\eqref{norm:determinant-sum}, evaluated at the three label vectors.
Consequently
\[
 \gamma_{k-1}\equiv
 \mathcal N(X(\xi_1),X(\xi_2),X(\xi_3))\pmod r,
\]
which is nonzero by \eqref{norm:distinct-labels}.
In particular $|\gamma_{k-1}|\ge1$ as an integer.

The integer
\[
 S_C=\sum_{u=0}^{k-1}\gamma_uB^u
\]
represents $\det C$ modulo $h$, because all omitted terms in $F(B)$
are divisible by $B^k$.
Since $B-1\ge100k^2L^3$, the coefficient bound gives
\begin{align}
 |S_C|&\le\frac{A_0}{B-1}(h-1)
           \le\frac3{50}(h-1)<\frac h2,
            \label{norm:centered}\\
 \left|\sum_{u=0}^{k-2}\gamma_uB^u\right|
       &\le\frac3{50}(B^{k-1}-1).
            \label{norm:lower-carry}
\end{align}
It follows that
\[
 |S_C|\ge B^{k-1}-\frac3{50}(B^{k-1}-1)
         =\frac{47B^{k-1}+3}{50}>\tau.
\]
By \eqref{norm:centered}, every other representative $S_C+mh$,
$m\in\mathbb Z\setminus\{0\}$, has absolute value greater than
$h/2>\tau$.  This proves the assertion.
\end{proof}

In particular, if $G\in\operatorname{SL}_3(R_h)$ and an integer matrix
$A$ satisfies $A\bmod h=GC$ and $|\det A|\le\tau$, then at least two
labels of $C$ coincide: multiplication by $G$ preserves its determinant
residue.  For three independently sampled labels, the probability of
such a coincidence is at most $3r^{-d}$ by the union bound.
